% ============================================================
%  Sistema de Plugins de OpenMarkov Core
%  Documento técnico
%
%  Compilar:
%    make              (genera PNGs con PlantUML y compila el PDF)
%  o manualmente:
%    pdflatex -interaction=nonstopmode plugins.tex   (dos pasadas)
% ============================================================
\documentclass[12pt,a4paper,spanish]{article}

\usepackage[utf8]{inputenc}
\usepackage[T1]{fontenc}
\usepackage[spanish]{babel}
\usepackage[margin=2.5cm]{geometry}
\usepackage{hyperref}
\usepackage{xcolor}
\usepackage{listings}
\usepackage{booktabs}
\usepackage{tabularx}
\usepackage{enumitem}
\usepackage{titlesec}
\usepackage{graphicx}
\usepackage{caption}
\usepackage{float}
\usepackage{fancyhdr}
\usepackage{microtype}
\usepackage{lmodern}

% ---- Colores ----
\definecolor{azulOM}{RGB}{30,90,160}
\definecolor{verdeManager}{RGB}{40,140,80}
\definecolor{grisFondo}{RGB}{245,245,245}

% ---- Hiperenlaces ----
\hypersetup{
    colorlinks=true,
    linkcolor=azulOM,
    urlcolor=azulOM,
    citecolor=azulOM,
    pdftitle={Sistema de Plugins de OpenMarkov Core},
    pdfauthor={OpenMarkov},
    pdfsubject={Arquitectura de extensibilidad}
}

% ---- Código fuente ----
\renewcommand{\lstlistingname}{Listado}

\lstdefinelanguage{JavaPlugin}{
    language=Java,
    morekeywords={@PotentialType,@FormatType,@Constraint,@NetworkTypeInfo,
                  @ImplementationRequirements,@RequiredConstructor,@RequiredMethod,
                  @AutoService,@Retention,@Target,RUNTIME,TYPE,
                  ConstraintBehavior,PotentialRole,PluginSearch,PluginLoader,
                  ExtensionTree,PotentialUtils,FormatManager,ConstraintManager,
                  NetworkTypeUtils,PNConstraint,NetworkType,Potential,
                  ProbNetReader,ProbNetWriter,InferenceAlgorithm,
                  PluginClassCategory,YES,NO,OPTIONAL,SelfClass},
    keywordstyle=\color{azulOM}\bfseries,
    commentstyle=\color{gray!70!black}\itshape,
    stringstyle=\color{verdeManager},
    basicstyle=\ttfamily\small,
    backgroundcolor=\color{grisFondo},
    frame=single,
    rulecolor=\color{gray!40},
    breaklines=true,
    numbers=left,
    numberstyle=\tiny\color{gray},
    numbersep=8pt,
    tabsize=4,
    showstringspaces=false,
    captionpos=b
}

% ---- Cabecera y pie ----
\pagestyle{fancy}
\fancyhf{}
\fancyhead[L]{\small\textcolor{azulOM}{\textbf{OpenMarkov Core}}}
\fancyhead[R]{\small Sistema de Plugins}
\fancyfoot[C]{\thepage}
\renewcommand{\headrulewidth}{0.4pt}

% ---- Secciones ----
\titleformat{\section}{\large\bfseries\color{azulOM}}{§\thesection}{1em}{}
\titleformat{\subsection}{\normalsize\bfseries\color{azulOM!80}}{\thesubsection}{1em}{}
\titleformat{\subsubsection}{\normalsize\bfseries\color{gray!60!black}}{\thesubsubsection}{1em}{}

% ---- Macro para figuras UML ----
\newcommand{\umlFig}[3]{%
  \begin{figure}[H]
    \centering
    \includegraphics[width=#2\textwidth]{#1.png}
    \caption{#3}
  \end{figure}%
}

% ============================================================
\begin{document}
% ============================================================

\begin{titlepage}
    \centering
    \vspace*{2cm}
    {\Huge\bfseries\color{azulOM} Sistema de Plugins\\[0.5em]
     \large\color{gray!80!black} \texttt{org.openmarkov.core}}\\[2cm]
    {\large Documentación técnica de la arquitectura de extensibilidad}\\[0.5em]
    {\normalsize Descubrimiento automático · Anotaciones · Managers · Contratos}\\[3cm]
    \rule{0.6\textwidth}{1pt}\\[1em]
    {\normalsize Módulo: \texttt{org.openmarkov.core}}\\
    {\normalsize Fecha: \today}
    \vfill
    \begin{center}
    \small\textit{Los diagramas UML se generan a partir de los archivos}\\
    \texttt{img/*.puml} \textit{mediante PlantUML (véase el Makefile adjunto).}
    \end{center}
\end{titlepage}

\tableofcontents
\clearpage

% ============================================================
\section{Introducción}
% ============================================================

En el contexto de OpenMarkov, un \textbf{plugin} es una clase Java que extiende o
implementa un tipo base previamente definido en el núcleo del sistema (por ejemplo,
un tipo de potencial, una restricción estructural o un formato de archivo), y que
el sistema descubre y registra \emph{automáticamente al arrancar}, sin que sea
necesario modificar ningún código existente. El desarrollador simplemente anota la
nueva clase con la anotación apropiada, la coloca en el classpath y el sistema la
incorpora como si siempre hubiera estado ahí.

Este mecanismo está diseñado para soportar la extensibilidad a largo plazo: nuevos
módulos de OpenMarkov pueden añadir tipos de red, algoritmos de inferencia o
formatos de archivo sin crear dependencias circulares ni modificar el núcleo.

\subsection{Cuatro pilares del sistema}
\label{sec:pilares}

El sistema de plugins se sustenta en cuatro mecanismos que trabajan de forma
coordinada. En conjunto, proporcionan un ciclo completo de extensión: el
desarrollador anota → el motor descubre → el manager registra → el dominio usa.

\begin{enumerate}[leftmargin=2em]
    \item \textbf{Anotaciones de descubrimiento} (\S\ref{sec:anotaciones}): marcan
          las clases como puntos de extensión y les asignan metadatos (nombre,
          comportamiento por defecto, versión, etc.). Son el contrato visible que
          el desarrollador de un plugin debe cumplir.
    \item \textbf{Motor de búsqueda de plugins} (\S\ref{sec:motor}): escanea el
          classpath al arrancar y construye un catálogo de todas las clases que
          cumplen los criterios de cada categoría. Opera de forma transparente,
          sin que el código de dominio tenga que invocarlo explícitamente.
    \item \textbf{Managers} (\S\ref{sec:managers}): singletons que consumen el
          catálogo construido por el motor y exponen al resto del sistema una API
          de alto nivel para obtener instancias de plugins por nombre, filtrarlos
          por compatibilidad o listar todos los disponibles.
    \item \textbf{Contrato de implementación} (\S\ref{sec:contrato}): la anotación
          \texttt{@ImplementationRequirements} documenta qué constructores y métodos
          deben existir en cada tipo de plugin. Los tests de integración verifican
          que ningún plugin rompa el contrato.
\end{enumerate}

\subsection{Cuatro categorías de plugin}
\label{sec:categorias}

Actualmente el núcleo de OpenMarkov define cuatro categorías de plugin, cada una
asociada a una anotación de descubrimiento y un manager específico. La tabla
siguiente resume su propósito y sus relaciones principales.

\medskip
\begin{tabularx}{\textwidth}{llXl}
\toprule
\textbf{Categoría} & \textbf{Anotación} & \textbf{Propósito} & \textbf{Manager} \\
\midrule
Potencial        & \texttt{@PotentialType}    & Distribuciones de probabilidad / utilidad & \texttt{PotentialUtils} \\
Restricción      & \texttt{@Constraint}        & Reglas estructurales de una red           & \texttt{ConstraintManager} \\
Tipo de red      & \texttt{@NetworkTypeInfo}   & Define el tipo de modelo probabilístico   & \texttt{NetworkTypeUtils} \\
Formato de arch. & \texttt{@FormatType}        & Lectura/escritura de redes en disco       & \texttt{FormatManager} \\
\bottomrule
\end{tabularx}

% ============================================================
\section{Vista de alto nivel}
\label{sec:vista-general}
% ============================================================

La figura~\ref{fig:general} muestra los cuatro grupos de actores del sistema y sus
relaciones principales: las anotaciones marcan los plugins, las clases base los
definen estructuralmente, los managers los registran y el motor de búsqueda
proporciona la infraestructura de descubrimiento. Los paquetes están ordenados
de izquierda a derecha según su papel en el ciclo de descubrimiento: las
anotaciones son el punto de entrada del desarrollador; las clases base definen
los contratos; los managers son los intermediarios que el dominio usa; y el motor
es el mecanismo interno que hace posible todo lo anterior.

\textbf{Leyenda de colores:}
\begin{itemize}
    \item \colorbox{yellow!40}{\strut Amarillo} — Anotaciones de descubrimiento
    \item \colorbox{orange!25}{\strut Naranja} — Clases base de los puntos de extensión
    \item \colorbox{green!30}{\strut Verde} — Managers (singletons de acceso)
    \item \colorbox{cyan!30}{\strut Cian} — Motor de búsqueda (\texttt{org.openmarkov.plugin})
    \item \colorbox{blue!15}{\strut Azul claro} — Implementaciones concretas de tipos de red
    \item \colorbox{red!12}{\strut Rosa} — Implementaciones concretas de formatos
    \item \colorbox{yellow!15}{\strut Crema} — Implementaciones concretas de potenciales/restricciones
\end{itemize}

\umlFig{img/fig01-general}{0.98}{Vista general del sistema de plugins. De izquierda a
  derecha: anotaciones de descubrimiento, clases base, managers y motor de búsqueda.}
\label{fig:general}

% ============================================================
\section{Motor de búsqueda de plugins}
\label{sec:motor}
% ============================================================

El paquete \texttt{org.openmarkov.plugin} proporciona la infraestructura de
descubrimiento. Es independiente del dominio y reutilizable en cualquier módulo
de OpenMarkov. Contiene tres clases principales y un enum.

\umlFig{img/fig02-motor}{0.92}{Motor de búsqueda: \texttt{PluginSearch}, \texttt{PluginLoader} y \texttt{ExtensionTree}}

\subsection{PluginClassCategory}

Enumera las fuentes de clases que el motor sabe escanear:

\begin{description}[leftmargin=2em]
    \item[\texttt{OPENMARKOV}] Clases cuyo nombre de paquete comienza por
          \texttt{org.openmarkov}. Es la categoría usada por todos los managers
          del núcleo. El escaneo se realiza mediante \textbf{ClassGraph} sobre el
          classpath completo.
    \item[\texttt{JAVA}] Clases del JDK, accedidas mediante el sistema de
          archivos virtual \texttt{jrt:/}.
    \item[\texttt{EXTERNAL\_DEPENDENCY}] Resto de clases del classpath (librerías
          de terceros, no OpenMarkov).
\end{description}

\subsection{PluginLoader}

Implementación interna (visibilidad de paquete). Responsabilidades:
\begin{enumerate}
    \item \textbf{Enumeración}: usa ClassGraph para obtener las rutas de todos los
          JARs y directorios del classpath y extrae los nombres de clase de cada
          archivo \texttt{.class}.
    \item \textbf{Carga en paralelo}: carga cada clase con
          \texttt{ClassLoader.getSystemClassLoader()} usando un stream paralelo.
    \item \textbf{Caché}: el resultado se guarda en \texttt{LOADED\_CLASSES}
          (un \texttt{HashMap} estático). Las categorías se inicializan de forma
          diferida y sincronizada: si dos threads solicitan la misma categoría
          simultáneamente, solo uno realiza el escaneo.
\end{enumerate}

\subsection{PluginSearch}

API fluida e inmutable construida sobre \texttt{Stream<Class<T>>}. Cada método
devuelve una nueva instancia de \texttt{PluginSearch}, permitiendo encadenar
operaciones sin efectos secundarios.

\begin{center}
\begin{tabular}{ll}
\toprule
\textbf{Método} & \textbf{Descripción} \\
\midrule
\texttt{init()} & Punto de entrada; obtiene el stream de \texttt{OPENMARKOV} \\
\texttt{full()} & Combina todas las categorías \\
\texttt{extending(Class<E>)} & Filtra subclases/implementaciones de \texttt{E} \\
\texttt{childrenOf(Class<E>)} & Como \texttt{extending} pero excluye la clase raíz \\
\texttt{annotatedWith(Class<A>)} & Filtra clases con la anotación \texttt{A} \\
\texttt{filter(Predicate)} & Filtro genérico \\
\texttt{stream()} & Operación terminal: devuelve el stream resultante \\
\texttt{list()} & Operación terminal: recoge en \texttt{List} \\
\texttt{extensionTree()} & Operación terminal: construye un \texttt{ExtensionTree} \\
\bottomrule
\end{tabular}
\end{center}

\begin{lstlisting}[language=JavaPlugin,
    caption={Ejemplo de PluginSearch encadenado para buscar restricciones}]
Stream<Class<? extends PNConstraint>> restricciones =
    PluginSearch.init()
        .annotatedWith(Constraint.class)
        .childrenOf(PNConstraint.class)
        .stream();
\end{lstlisting}

\subsection{ExtensionTree}

Estructura de árbol que refleja la jerarquía de herencia. Se construye a partir
de una clase raíz y un \texttt{Stream} de sus subclases. El método
\texttt{breathFirstLevelOrderQueue()} recorre el árbol en \textbf{anchura (BFS)},
garantizando que las clases más cercanas a la raíz se procesen antes que las
hojas. Esto es crucial para \texttt{PotentialUtils}: cuando una clase abstracta
lleva \texttt{@PotentialType} y varias clases hoja la heredan, la clase más
cercana a la raíz prevalece en el mapa de nombres.

% ============================================================
\section{Anotaciones de descubrimiento y de contrato}
\label{sec:anotaciones}
% ============================================================

Las anotaciones del sistema se dividen en dos grupos: las de \emph{descubrimiento}
(que marcan plugins) y las de \emph{contrato} (que documentan requisitos de
implementación).

Todas las anotaciones de descubrimiento comparten dos meta-anotaciones de Java:

\begin{itemize}
    \item \texttt{@Retention(RetentionPolicy.RUNTIME)}: indica que la anotación
          debe estar disponible en tiempo de ejecución (y no solo en tiempo de
          compilación). Esto permite que \texttt{PluginSearch} la lea mediante
          reflexión durante el escaneo del classpath.
    \item \texttt{@Target(ElementType.TYPE)}: indica que la anotación solo puede
          aplicarse a declaraciones de tipo (clases, interfaces, enumeraciones y
          anotaciones). No puede aplicarse a métodos, campos ni parámetros.
\end{itemize}

\umlFig{img/fig03-anotaciones}{0.90}{Todas las anotaciones del sistema de plugins}

\subsection{Anotaciones de descubrimiento}

\subsubsection{\texttt{@PotentialType}}

Marca las subclases de \texttt{Potential} que deben ser descubiertas. El atributo
\texttt{names} acepta un array para permitir múltiples aliases del mismo tipo,
útil para mantener compatibilidad con nombres históricos de archivos XML:

\begin{lstlisting}[language=JavaPlugin, caption={Uso de @PotentialType}]
@PotentialType(names = {"Table", "ProbTable"})
public class TablePotential extends AbstractIndexedPotential { ... }

@PotentialType(names = "Uniform")
public class UniformPotential extends Potential { ... }

@PotentialType(names = "ICIModel")  // heredada por MaxPotential y MinPotential
public abstract class ICIPotential extends Potential { ... }
\end{lstlisting}

\subsubsection{\texttt{@Constraint}}

Marca las subclases de \texttt{PNConstraint}. El atributo \texttt{defaultBehavior}
determina el comportamiento global de la restricción antes de que los tipos de
red lo sobrescriban:

\begin{description}[leftmargin=2em]
    \item[\texttt{YES}] Activa en todos los tipos de red por defecto.
    \item[\texttt{NO}] Inactiva por defecto; un tipo de red puede activarla.
    \item[\texttt{OPTIONAL}] Inactiva por defecto, pero puede activarla el usuario.
\end{description}

\subsubsection{\texttt{@NetworkTypeInfo}}

Proporciona el nombre técnico (usado en ficheros XML), el nombre visual (mostrado
en la GUI) y nombres alternativos para compatibilidad con versiones anteriores.

\subsubsection{\texttt{@FormatType}}

Incluye la extensión de archivo y la versión del formato para la resolución
automática del lector/escritor. Permite registrar múltiples versiones del mismo
formato (por ejemplo, \texttt{pgmx~0.1} y \texttt{pgmx~0.2}).

\subsection{Anotaciones de contrato}
\label{sec:contrato}

\subsubsection{\texttt{@ImplementationRequirements}}

Documenta qué constructores y métodos deben existir en \emph{todas} las subclases
concretas. \textbf{No impone restricciones en tiempo de ejecución}: la verificación
se realiza en los tests de integración.

\begin{lstlisting}[language=JavaPlugin,
    caption={@ImplementationRequirements en Potential}]
@ImplementationRequirements(
    requiresOneOfTheseConstructors = {
        @RequiredConstructor({List.class, CycleLength.class}),   // DBN
        @RequiredConstructor({List.class, PotentialRole.class}), // estándar
        @RequiredConstructor(List.class),                        // mínimo
        @RequiredConstructor(SelfClass.class)                    // copy ctor
    },
    requiresMethods = @RequiredMethod(
        methodKind = RequiredMethod.MethodKind.Static,
        methodName = "validate",
        returnType = Boolean.class,
        parameters = {Node.class, List.class, PotentialRole.class}
    )
)
public abstract class Potential implements Localizable { ... }
\end{lstlisting}

La clase marcador \texttt{SelfClass} representa «el tipo de la propia subclase»,
exigiendo un constructor copia de la forma \texttt{SubTipo(SubTipo other)}.

Para \texttt{PNConstraint} el requisito es más sencillo: solo exige un constructor
sin argumentos, necesario para que \texttt{ConstraintManager} las instancie por
reflexión:

\begin{lstlisting}[language=JavaPlugin,
    caption={@ImplementationRequirements en PNConstraint}]
@ImplementationRequirements(
    requiresOneOfTheseConstructors = @RequiredConstructor({}) // sin args
)
public abstract class PNConstraint { ... }
\end{lstlisting}

% ============================================================
\section{Managers}
\label{sec:managers}
% ============================================================

Un \textbf{manager} es un singleton que actúa como punto de acceso centralizado
a una categoría de plugins. Sus responsabilidades son tres: (1)~descubrir todas
las clases de plugins disponibles al arrancar, usando \texttt{PluginSearch};
(2)~mantener un registro interno (mapa de nombres a clases o lista de instancias)
para que las búsquedas posteriores sean eficientes; y (3)~exponer una API de alto
nivel que el código de dominio puede usar sin conocer los detalles del sistema de
reflexión subyacente.

Todos los managers se inicializan de forma diferida la primera vez que se accede
a ellos, y el resultado del descubrimiento se cachea de forma permanente durante
la vida de la JVM.

\umlFig{img/fig04-managers}{0.98}{Los cuatro managers (columna central) y sus relaciones
  con el motor de búsqueda (izquierda) y las anotaciones que leen (derecha).}

\subsection{PotentialUtils}
\label{sec:potentialutils}

\texttt{PotentialUtils} es el manager de potenciales. Su responsabilidad es
mantener el catálogo completo de tipos de potencial disponibles, permitir su
búsqueda por nombre, instanciarlos de forma segura mediante reflexión, y
determinar qué tipos son compatibles con un nodo concreto. Es la clase que
permite que el sistema (y la GUI) traten los potenciales como entidades
intercambiables sin conocer sus clases concretas.

El mapa \texttt{POTENTIALS\_BY\_NAME} se construye en el inicializador estático
recorriendo el árbol de extensión de \texttt{Potential} en amplitud (BFS).
Esto garantiza que, si hay un conflicto de nombre entre clase padre e hija, la
más cercana a la raíz prevalece.

El método \texttt{instanciateSafely} prueba constructores en este orden:
\begin{enumerate}
    \item \texttt{(List<Variable>, CycleLength)} — para redes dinámicas (DBN).
    \item \texttt{(List<Variable>, PotentialRole)} — constructor estándar.
    \item \texttt{(List<Variable>)} — constructor mínimo.
\end{enumerate}

El método \texttt{getFilteredPotentialClasses(Node)} usa reflexión para invocar
el método estático \texttt{validate(Node, List<Variable>, PotentialRole)} de cada
potencial y devuelve solo los compatibles con el nodo dado.

\subsection{ConstraintManager}
\label{sec:constraintmanager}

\texttt{ConstraintManager} es el manager de restricciones. Su responsabilidad es
conocer el conjunto completo de restricciones disponibles en el sistema y, para
cada tipo de red concreto, construir la lista activa de restricciones que deben
verificarse. Actúa como mediador entre el comportamiento global por defecto de
cada restricción y las personalizaciones que cada tipo de red puede introducir.

La lógica de \texttt{buildConstraintList} opera en dos fases:
\begin{enumerate}
    \item Recoge todas las restricciones con \texttt{defaultBehavior == YES}
          (y también las \texttt{OPTIONAL} si el parámetro \texttt{includeOptionals}
          es \texttt{true}).
    \item Aplica las \emph{sobrescrituras} del tipo de red concreto
          (\texttt{NetworkType.getOverwrittenConstraints()}): si una restricción
          está marcada \texttt{YES} en el tipo, se añade; si está marcada
          \texttt{NO}, se elimina de la lista.
\end{enumerate}

\begin{lstlisting}[language=JavaPlugin,
    caption={Lógica simplificada de ConstraintManager.buildConstraintList}]
public List<PNConstraint> buildConstraintList(NetworkType type, boolean optionals) {
    List<PNConstraint> result = new ArrayList<>();
    // Fase 1: comportamientos por defecto globales
    for (var entry : defaultConstraintBehaviors.entrySet()) {
        if (entry.getValue() == YES || (optionals && entry.getValue() == OPTIONAL))
            result.add(instantiate(entry.getKey()));
    }
    // Fase 2: sobrescrituras del tipo de red
    for (var entry : type.getOverwrittenConstraints().entrySet()) {
        if (entry.getValue() == YES)
            result.add(instantiate(entry.getKey()));
        else if (entry.getValue() == NO)
            result.removeIf(c -> c.getClass() == entry.getKey());
    }
    return result;
}
\end{lstlisting}

\subsection{NetworkTypeUtils}
\label{sec:networktypeutils}

\texttt{NetworkTypeUtils} es el manager de tipos de red. Su responsabilidad es
mantener el catálogo de todos los tipos de red disponibles en el sistema y
proporcionar acceso a ellos por nombre. Es el componente que permite a los parsers
de archivos y al código de dominio obtener la instancia singleton correcta de un
tipo de red a partir de su nombre textual (por ejemplo, al cargar un archivo XML
que declara \texttt{type="BayesianNetwork"}).

Todos los tipos de red son singletons. \texttt{NetworkTypeUtils.safeInstanciate}
los obtiene llamando al método estático \texttt{getUniqueInstance()} por reflexión,
evitando la instanciación duplicada.

La búsqueda por nombre prueba primero el atributo \texttt{name} de la anotación
y, si no hay coincidencia, prueba los \texttt{alternativeNames}. Esto garantiza
compatibilidad con archivos guardados con versiones anteriores de OpenMarkov.

\subsection{FormatManager}
\label{sec:formatmanager}

\texttt{FormatManager} es el manager de formatos de archivo. Su responsabilidad
es descubrir todos los lectores (\texttt{ProbNetReader}) y escritores
(\texttt{ProbNetWriter}) disponibles, instanciarlos al arrancar y proporcionar el
reader o writer correcto para cualquier archivo que se quiera abrir o guardar.

A diferencia de los otros managers, \texttt{FormatManager} \textbf{instancia
inmediatamente} todos los plugins descubiertos en su constructor privado. El motivo
es que la operación de abrir un archivo debe ser lo más rápida posible; pagar el
coste de instanciación al arranque es preferible a pagarlo en cada apertura.

La resolución del reader para un archivo concreto combina extensión y versión:
\begin{enumerate}
    \item Extrae la extensión del nombre de archivo.
    \item Para formatos distintos de \texttt{.elv}, lee la versión del propio
          contenido XML del archivo.
    \item Busca el reader cuya anotación \texttt{@FormatType} coincide con ambos
          valores (usando \texttt{startsWith} para la versión).
\end{enumerate}

% ============================================================
\section{Subsistemas de plugins}
\label{sec:subsistemas}
% ============================================================

En esta sección se documenta cada categoría de plugin en detalle: las clases base
que definen su contrato, las implementaciones concretas más relevantes y el
comportamiento específico de su sistema de descubrimiento. Las cinco subsecciones
siguientes corresponden a las cuatro categorías de plugin más las clases base que
las sustentan.

% ----------------------------------------------------------
\subsection{Clases base de los puntos de extensión}
\label{sec:bases}
% ----------------------------------------------------------

Las clases base definen el contrato que todo plugin debe cumplir: qué métodos
deben implementarse, qué constructor se exige y qué comportamiento estático
se espera. Son las únicas clases del núcleo que conocen los tipos de plugin de
forma abstracta; el resto del sistema opera siempre a través de los managers,
nunca directamente sobre las implementaciones concretas.

\umlFig{img/fig05-bases}{0.85}{Jerarquía de las clases base que actúan como puntos de extensión}

La tabla siguiente resume las obligaciones de cada clase base:

\medskip
\begin{tabularx}{\textwidth}{lX}
\toprule
\textbf{Clase base} & \textbf{Contrato para subclases} \\
\midrule
\texttt{Potential}      & 7 métodos abstractos + \texttt{static validate()} \\
\texttt{PNConstraint}   & \texttt{checkProbNet()} + constructor sin args \\
\texttt{NetworkType}    & \texttt{static getUniqueInstance()} + \texttt{@NetworkTypeInfo} \\
\texttt{ProbNetReader}  & \texttt{loadProbNetInfo()} + \texttt{loadProbNet()} + ctor sin args \\
\texttt{ProbNetWriter}  & \texttt{writeProbNet()} (2 sobrecargas) + ctor sin args \\
\bottomrule
\end{tabularx}

% ----------------------------------------------------------
\subsection{Subsistema de Potenciales}
\label{sec:potenciales}
% ----------------------------------------------------------

Los potenciales son el tipo de plugin más rico del sistema. Representan cualquier
función matemática que mapea configuraciones de variables a valores numéricos:
tablas de probabilidad condicional (CPTs), distribuciones paramétricas, formas
canónicas como noisy-OR, árboles de decisión, etc. La gran variedad de tipos
disponibles y la necesidad de filtrarlos por compatibilidad con el nodo hace
de \texttt{PotentialUtils} el manager más complejo del sistema.

\umlFig{img/fig06-potenciales}{0.88}{Subsistema de Potenciales: anotación, manager y jerarquía de implementaciones}

\subsubsection{Descubrimiento y herencia de la anotación}

El inicializador estático de \texttt{PotentialUtils} usa
\texttt{.extending(Potential.class)} en lugar de
\texttt{.annotatedWith(@PotentialType)}. Este detalle es deliberado: permite que
una clase abstracta intermedia (como \texttt{ICIPotential}) lleve la anotación
\texttt{@PotentialType(names = "ICIModel")} y sus hijos concretos
(\texttt{MaxPotential}, \texttt{MinPotential}) la \emph{hereden} sin necesidad de
repetirla.

\begin{lstlisting}[language=JavaPlugin,
    caption={Descubrimiento de potenciales en el inicializador estático de PotentialUtils}]
static {
    var byLevel = PluginSearch.init()
        .extending(Potential.class)   // hereda la anotación de clases padre
        .extensionTree()
        .breathFirstLevelOrderQueue()
        .stream()
        .map(ExtensionTree::getCurrentClass)
        .toList();

    byLevel.forEach(cls ->
        PotentialUtils.getNames(cls).forEach(name ->
            potentialsByName.put(name, cls)));   // registra todos los aliases
    POTENTIALS_BY_NAME = potentialsByName;
}
\end{lstlisting}

\subsubsection{Filtrado por compatibilidad}

Cuando la GUI necesita saber qué potenciales son válidos para un nodo concreto,
invoca \texttt{getFilteredPotentialClasses(node)}. Este método llama por reflexión
al método estático \texttt{validate(Node, List<Variable>, PotentialRole)} de cada
potencial. La implementación por defecto en \texttt{Potential} devuelve
\texttt{true}; las subclases la sobrescriben. Por ejemplo,
\texttt{TablePotential.validate} rechaza variables de tipo \texttt{NUMERIC}.

\subsubsection{Indexación interna de TablePotential}

\texttt{TablePotential} almacena la tabla en un array plano \texttt{double[] values}.
La posición de una configuración $(x_0, x_1, \ldots)$ es:
\[
\mathrm{pos} = x_0 \cdot \mathrm{offsets}[0] + x_1 \cdot \mathrm{offsets}[1] + \cdots
\]
donde \texttt{offsets[0]=1} y \texttt{offsets[i] = dimensions[i-1] * offsets[i-1]}.
Esta indexación lexicográfica permite un acceso eficiente por indexación directa.

% ----------------------------------------------------------
\subsection{Subsistema de Restricciones}
\label{sec:restricciones}
% ----------------------------------------------------------

Las restricciones son reglas estructurales que una red debe cumplir en todo
momento. Se verifican antes de ejecutar cada edición (\S\texttt{PNEdit}) para
garantizar que ninguna operación deja la red en un estado inválido. Son instancias
sin estado (o con estado mínimo inmutable) que el sistema puede crear, aplicar y
descartar sin efectos secundarios. El método \texttt{checkProbNet} añade
excepciones al \texttt{ConstraintChecker} en lugar de lanzarlas directamente,
lo que permite acumular varias violaciones y reportarlas juntas en un único error.

\umlFig{img/fig07-restricciones}{0.92}{Subsistema de Restricciones: anotación, manager e implementaciones en dos columnas}

\subsubsection{Interacción con el sistema de ediciones}

\texttt{PNConstraint} implementa \texttt{PNEditListener}. Antes de ejecutar una
edición, esta llama a \texttt{checkConstraintsWillBeMet(checker)}, que añade al
\texttt{ConstraintChecker} solo las violaciones relevantes a la operación que va
a realizar, sin tener que reverificar toda la red desde cero.

% ----------------------------------------------------------
\subsection{Subsistema de Tipos de Red}
\label{sec:tiposred}
% ----------------------------------------------------------

Los tipos de red son el mecanismo por el cual OpenMarkov soporta distintos
formalismos probabilísticos (redes bayesianas, diagramas de influencia, MDPs,
etc.) con un único núcleo común. Cada tipo de red personaliza el conjunto de
restricciones activas, determinando qué operaciones son válidas sobre ella. Son
singletons inmutables porque su función es puramente declarativa: no almacenan
datos de la red, solo sus reglas de validez.

\umlFig{img/fig08-tiposred}{0.90}{Subsistema de Tipos de Red: anotación, manager e implementaciones}

\begin{lstlisting}[language=JavaPlugin,
    caption={Estructura típica de un tipo de red (BayesianNetworkType)}]
@NetworkTypeInfo(name = "BayesianNetwork", visualName = "Bayesian Network")
public final class BayesianNetworkType extends NetworkType {
    private static final BayesianNetworkType INSTANCE = new BayesianNetworkType();

    private BayesianNetworkType() {
        super();
        // Activa una restricción que por defecto global es NO
        overrideConstraintBehavior(OnlyChanceNodes.class, ConstraintBehavior.YES);
    }

    public static synchronized BayesianNetworkType getUniqueInstance() {
        return INSTANCE;
    }
}
\end{lstlisting}

% ----------------------------------------------------------
\subsection{Subsistema de Formatos de Archivo}
\label{sec:formatos}
% ----------------------------------------------------------

Los plugins de formato son la puerta de entrada y salida del sistema: permiten
leer redes probabilísticas desde archivos en disco y guardarlas de vuelta. Están
diseñados para que añadir soporte a un nuevo formato (o una nueva versión de uno
existente) no requiera modificar ningún código del núcleo; basta con añadir
una nueva clase anotada con \texttt{@FormatType}.

A diferencia de los otros subsistemas, los plugins de formato se instancian de
forma \emph{inmediata} al arrancar: el \texttt{FormatManager} crea una instancia
de cada reader y cada writer en su constructor privado. Esto garantiza que
cualquier error de configuración (clase no encontrada, constructor ausente) se
detecte durante el arranque y no al intentar abrir un archivo.

Un mismo módulo puede registrar tanto un reader como un writer para el mismo
formato (extensión y versión) anotando clases separadas. El
\texttt{FormatManager} los mantiene en listas distintas y los resuelve
independientemente.

\umlFig{img/fig09-formatos}{0.88}{Subsistema de Formatos de Archivo: anotación, manager, interfaces y ejemplos de implementación}

La resolución del reader para un archivo concreto combina tres datos: extensión
del archivo, versión leída del contenido XML y la lista de readers disponibles.
El método \texttt{getProbNetReaderInstanceFor(ext, version)} busca el primero
cuya anotación \texttt{@FormatType} coincida con ambos valores, usando
\texttt{startsWith} para la versión (por ejemplo, \texttt{"0"} coincide con
\texttt{"0.2"}).

% ============================================================
\section{Flujos de descubrimiento: diagramas de secuencia}
\label{sec:secuencia}
% ============================================================

Los diagramas de secuencia siguientes muestran el flujo completo de descubrimiento
para tres categorías de plugins: potenciales (el más complejo, con árbol de herencia
y mapa de nombres), restricciones (con la doble fase de comportamientos globales y
sobrescrituras por tipo de red) y formatos (con instanciación inmediata al arranque).
En los tres casos, el flujo sigue el mismo patrón general: la primera referencia al
manager dispara el descubrimiento vía \texttt{PluginSearch} y \texttt{PluginLoader};
el resultado se cachea; las llamadas posteriores son simples consultas al caché.

\subsection{Descubrimiento de potenciales}

El diagrama de la figura~\ref{fig:seq-potenciales} muestra el flujo completo desde
que la JVM referencia por primera vez a \texttt{PotentialUtils} (disparando el
inicializador estático) hasta que el código externo obtiene una instancia de
potencial. Hay dos fases bien diferenciadas: la \emph{fase de construcción del
catálogo} (carga del classpath, filtrado, BFS del árbol de herencia y registro
en el mapa) y la \emph{fase de uso} (búsqueda en el mapa e instanciación por
reflexión). La fase de uso es $O(1)$ gracias al mapa preconstruido.

Un punto clave del flujo es que \texttt{PluginSearch} usa
\texttt{.extending(Potential.class)} y no \texttt{.annotatedWith(@PotentialType)},
lo que garantiza que las clases que heredan la anotación de una clase padre sean
también descubiertas.

\umlFig{img/fig10-seq-potenciales}{0.95}{Secuencia de descubrimiento e instanciación de potenciales}
\label{fig:seq-potenciales}

\subsection{Construcción de la lista de restricciones}

El diagrama de la figura~\ref{fig:seq-restricciones} muestra cómo la construcción
de un \texttt{ProbNet} con tipo \texttt{BayesianNetworkType} resulta en una lista
de restricciones activas. El flujo tiene dos momentos temporalmente separados: el
\emph{arranque de la JVM}, en el que \texttt{ConstraintManager} se inicializa y
descubre todas las restricciones disponibles; y la \emph{creación de una red}
concreta, en la que \texttt{buildConstraintList} aplica los comportamientos por
defecto y luego las sobrescrituras del tipo de red.

El efecto de la doble fase se ve claramente con \texttt{OnlyChanceNodes}: su
\texttt{defaultBehavior} global es \texttt{NO}, por lo que no entra en la lista
inicial; pero \texttt{BayesianNetworkType} la sobrescribe con \texttt{YES}, de
modo que se añade en la segunda fase.

\umlFig{img/fig11-seq-restricciones}{0.93}{Secuencia de construcción de la lista de restricciones para una red bayesiana}
\label{fig:seq-restricciones}

\subsection{Resolución de formato al abrir un archivo}

El diagrama de la figura~\ref{fig:seq-formatos} ilustra el flujo del
\texttt{FormatManager}. A diferencia de los otros dos, aquí el arranque es más
costoso: el manager instancia inmediatamente \emph{todos} los readers y writers
disponibles. La recompensa llega en el momento de abrir un archivo: la resolución
del reader es una simple búsqueda sobre una lista de instancias ya creadas,
sin reflexión ni instanciación en ese momento.

El diagrama también muestra la lectura de la versión del XML como paso intermedio
entre la extracción de la extensión y la búsqueda del reader. Este paso es
necesario porque el mismo formato puede tener múltiples versiones con parsers
distintos (por ejemplo, \texttt{pgmx~0.1} y \texttt{pgmx~0.2}).

\umlFig{img/fig12-seq-formatos}{0.95}{Secuencia de resolución y uso del lector de formato}
\label{fig:seq-formatos}

% ============================================================
\section{Cómo añadir un nuevo plugin}
\label{sec:extension}
% ============================================================

El sistema está diseñado para que añadir un nuevo plugin requiera únicamente:
\begin{enumerate}
    \item Crear la clase en cualquier módulo cuyo JAR esté en el classpath.
    \item Anotar con la anotación apropiada.
    \item Cumplir el contrato de \texttt{@ImplementationRequirements}.
\end{enumerate}
No es necesario modificar ninguna clase existente del núcleo.

\subsection{Nuevo Potencial}

\begin{lstlisting}[language=JavaPlugin, caption={Implementación de un nuevo potencial}]
@PotentialType(names = {"MiPotencial", "MyPotential"})
public class MiPotencial extends Potential {

    // Constructor estándar (requerido por @ImplementationRequirements)
    public MiPotencial(List<Variable> variables, PotentialRole role) {
        super(variables, role);
    }

    // Constructor copia (requerido: SelfClass en @ImplementationRequirements)
    public MiPotencial(MiPotencial other) { super(other); }

    // Filtro de compatibilidad (requerido: @RequiredMethod validate)
    public static boolean validate(Node node, List<Variable> variables,
                                   PotentialRole role) {
        return variables.stream()
            .allMatch(v -> v.getVariableType() == VariableType.FINITE_STATES);
    }

    @Override
    public TablePotential tableProject(EvidenceCase evidence,
                                       InferenceOptions options,
                                       List<TablePotential> projected) { ... }
    @Override public Potential project(EvidenceCase evidence) { ... }
    @Override public Potential copy()  { return new MiPotencial(this); }
    @Override public boolean isUncertain() { return false; }
    @Override public void scalePotential(double scale) { ... }
    @Override public Potential reorder(List<Variable> newOrder) { ... }
    @Override public Potential reorder(Variable v, State[] newOrder) { ... }
}
// Acceso: PotentialUtils.getClassByName("MiPotencial")
\end{lstlisting}

\subsection{Nueva Restricción}

\begin{lstlisting}[language=JavaPlugin, caption={Implementación de una nueva restricción}]
@Constraint(name = "MiRestriccion", defaultBehavior = ConstraintBehavior.YES)
public class MiRestriccion extends PNConstraint {

    // Constructor sin argumentos (requerido por @ImplementationRequirements)
    public MiRestriccion() { super(); }

    @Override
    public void checkProbNet(ProbNet probNet, ConstraintChecker checker) {
        for (Node node : probNet.getNodes()) {
            if (condicionDeViolacion(node))
                checker.addException(new ConstraintViolatedException.MiError(this, node));
        }
    }
}
// Un tipo de red puede desactivarla con:
//   overrideConstraintBehavior(MiRestriccion.class, ConstraintBehavior.NO);
\end{lstlisting}

\subsection{Nuevo Tipo de Red}

\begin{lstlisting}[language=JavaPlugin, caption={Implementación de un nuevo tipo de red}]
@NetworkTypeInfo(
    name        = "MiRed",
    visualName  = "Mi Tipo de Red",
    alternativeNames = {"MyNet"}
)
public final class MiRedType extends NetworkType {

    private static final MiRedType INSTANCE = new MiRedType();

    private MiRedType() {
        super();
        overrideConstraintBehavior(OnlyChanceNodes.class, ConstraintBehavior.YES);
        overrideConstraintBehavior(NoBackwardLink.class, ConstraintBehavior.NO);
    }

    // Obligatorio: NetworkTypeUtils lo invoca por reflexión
    public static synchronized MiRedType getUniqueInstance() { return INSTANCE; }
}
\end{lstlisting}

\subsection{Nuevo Formato de Archivo}

\begin{lstlisting}[language=JavaPlugin, caption={Implementación de un nuevo formato de archivo}]
@FormatType(
    name        = "MiFormatoReader",
    version     = "1.0",
    extension   = "mif",
    description = "MiFormato.description"   // clave i18n para la GUI
)
public class MiFormatoReader implements ProbNetReader {

    public MiFormatoReader() { }    // constructor sin args obligatorio

    @Override
    public ProbNetInfo loadProbNetInfo(String netName, InputStream file)
            throws IOException, ParserException {
        return new ProbNetInfo(...);    // metadatos ligeros
    }

    @Override
    public ProbNet loadProbNet(String netName, InputStream file)
            throws IOException, ParserException {
        ProbNet net = new ProbNet(MiRedType.getUniqueInstance());
        // parseo completo...
        return net;
    }
}
// Acceso: FormatManager.getInstance().getProbNetReader(urlArchivoMif)
\end{lstlisting}

% ============================================================
\section{Mapa completo del sistema}
\label{sec:mapa}
% ============================================================

La figura~\ref{fig:mapa} muestra cómo todos los componentes se ensamblan en
una vista única. El diagrama está organizado en dos franjas horizontales: en la
superior aparecen el motor de búsqueda, los managers y el código de dominio
(de izquierda a derecha según su papel en el flujo de descubrimiento); en la
inferior aparece el contenido del classpath con todos los tipos de plugins,
dispuestos horizontalmente para mostrar que son entidades independientes que
conviven en el mismo espacio de clases.

Las flechas del motor al classpath reflejan la carga por \texttt{PluginLoader};
las flechas de los managers al motor reflejan las búsquedas por \texttt{PluginSearch};
y las flechas del dominio a los managers reflejan el uso en tiempo de ejecución
(crear redes, validar ediciones, instanciar potenciales).

\umlFig{img/fig13-mapa}{0.95}{Mapa completo del sistema de plugins}
\label{fig:mapa}

% ============================================================
\section{Consideraciones de rendimiento y ciclo de vida}
\label{sec:rendimiento}
% ============================================================

\subsection{Coste de inicialización}

El escaneo del classpath por \texttt{PluginLoader} es la operación más costosa.
Para mitigarlo:
\begin{itemize}
    \item El escaneo se realiza \textbf{una sola vez} por categoría y el resultado
          se cachea en \texttt{LOADED\_CLASSES} (mapa estático final).
    \item La carga de clases individuales se hace \textbf{en paralelo} (stream
          paralelo).
    \item Los mapas de nombres se construyen una sola vez en inicializadores
          estáticos.
    \item \texttt{FormatManager} paga el coste de instanciación al arranque para
          que las aperturas de archivo posteriores sean inmediatas.
\end{itemize}

\subsection{Thread safety}

\begin{itemize}
    \item \texttt{PluginLoader.initializeCategory} está sincronizado.
    \item \texttt{ConstraintManager} y \texttt{FormatManager} son singletons con
          campo \texttt{static final} (garantía de la JVM).
    \item Los tipos de red tienen \texttt{getUniqueInstance()} sincronizado.
\end{itemize}

\subsection{Extensibilidad en tiempo de ejecución}

El sistema \emph{no} soporta recarga dinámica de plugins en caliente. Para añadir
un plugin es necesario reiniciar la JVM con el nuevo JAR en el classpath.

\subsection{Resumen comparativo}

\bigskip
\begin{tabularx}{\textwidth}{lXlll}
\toprule
\textbf{Manager} & \textbf{Búsqueda PluginSearch} & \textbf{Init} & \textbf{Instanciación} & \textbf{Caché} \\
\midrule
\texttt{PotentialUtils}    & \texttt{.extending(Potential)}             & static init & Diferida por nombre & \texttt{HashMap} \\
\texttt{ConstraintManager} & \texttt{.annotatedWith(@Constraint)}\newline\texttt{.childrenOf(PNConstraint)} & ctor privado & Al construir lista & \texttt{HashMap} \\
\texttt{NetworkTypeUtils}  & \texttt{.annotatedWith(@NetworkTypeInfo)}\newline\texttt{.childrenOf(NetworkType)} & static init & \texttt{getUniqueInstance()} & \texttt{List} \\
\texttt{FormatManager}     & \texttt{.annotatedWith(@FormatType)}\newline\texttt{.extending(Reader/Writer)} & ctor privado & Inmediata & \texttt{List} \\
\bottomrule
\end{tabularx}

\end{document}
